\section{指令扩展与硬件加速}\label{sec:extensions}

\subsection{基础指令能力}

处理器实现 RV32I 整数指令集、RV32M 乘除法指令、Zicsr CSR 访问指令以及机器态异常返回，
并实现了部分 RISC-V 标量 B 扩展指令。表~\ref{tab:implemented-extensions} 按标准扩展名称列出实际
实现范围；这里的指令列表即为本设计支持的子集，不表示已经覆盖对应扩展的全部指令。

\begin{table}[H]\centering\scriptsize
\caption{提交 \texttt{8591368} 的标准扩展能力与已实现指令}\label{tab:implemented-extensions}
\begin{tabularx}{\linewidth}{@{}C{0.095\linewidth}p{0.57\linewidth}Y@{}}
\toprule
\textbf{扩展}&\textbf{已实现的指令或能力}&\textbf{实现通路}\\\midrule
RV32M&\code{mul/mulh/mulhsu/mulhu/div/divu/rem/remu}&DSP 流水乘法与迭代除法\\
Zicsr&\code{csrrw/csrrs/csrrc} 及三种立即数形式&CSR 读改写通路\\
Zba&\code{sh1add/sh2add/sh3add}&单周期移位加法\\
Zbb&\code{clz/ctz/sext.b/sext.h/min/minu/max/maxu/rol/andn/orn/xnor/rev8}&短组合或 32 轮计数通路\\
Zbs&\code{bext/bexti/bset/bseti/bclr/bclri/binv/binvi}&单周期位操作\\
Zbkb&\code{pack}&单周期半字拼接\\
\bottomrule
\end{tabularx}
\normalsize\end{table}

\subsection{乘除法与 B 扩展多周期控制}

乘法器按操作类型选择普通低位、全宽高位或窄操作数快速通路；除法器采用恢复除法，每拍执行移位、
比较与减法，最后单独完成符号修正。当前 RTL 的迭代 B 单元只保留 \code{clz/ctz}，逐位执行 32 轮；
移位加法、符号扩展、最小/最大值、逻辑否定组合、\code{rol} 左旋、字节反序和单比特操作保留在短组合
通路。当前支持范围仅以表~\ref{tab:implemented-extensions} 为准；旧迭代单元中的其余操作不再支持。
\code{pack} 由当前单周期快速整数通路实现。

多周期请求被接受后，流水线相关信息仍保留目的寄存器号，但将 \code{dataValid} 保持为 0；因此依赖
该结果的年轻指令会在 IDU 停顿，而无关指令不会把未完成的数据误认为可前递结果。执行单元寄存最终
结果后置位 \code{dataValid}，待后级 \code{ready} 有效时完成握手并释放占用。该协议统一覆盖乘法、
除法、迭代 B 操作以及后述阻塞式加速单元。

\begin{figure}[H]\centering
\begin{tikzpicture}[s/.style={draw,rounded corners=1.2pt,minimum width=3.1cm,minimum height=0.82cm,align=center,fill=blue!6},d/.style={draw,diamond,aspect=2,align=center,fill=orange!7},a/.style={-Stealth,thick},font=\small]
\node[s] (idle) {Idle：接受多周期指令}; \node[s,below=0.7cm of idle] (busy) {Busy：保持 rd 占用\\dataValid=0};
\node[d,below=0.7cm of busy] (valid) {结果有效？}; \node[d,below=0.8cm of valid] (ready) {后级 ready？};
\node[s,below left=0.8cm and 1.4cm of ready] (hold) {Done：寄存结果}; \node[s,below right=0.8cm and 1.4cm of ready] (emit) {完成握手并释放};
\draw[a](idle)--node[right]{in.fire}(busy);\draw[a](busy)--(valid);\draw[a](valid.west)--++(-1,0)|-node[left]{否}(busy.west);
\draw[a](valid)--node[right]{是}(ready);\draw[a](ready)--node[above]{否}(hold);\draw[a](ready)--node[above]{是}(emit);\draw[a](hold)--(emit);
\end{tikzpicture}\caption{多周期执行单元握手与相关保持}\end{figure}


\subsection{B 扩展的动态热点与性能边界}

本小节记录早期 B 扩展探索的动态剖析，用于解释后续为何将更多热点移入专用硬件单元；其中
\code{clmul/clmulh} 等操作不属于提交 \code{8591368} 当前保留的硬件子集。为区分编译器利用 ISA 后的
指令数变化与具体硬件实现的周期代价，本文采用相同 CoreMark 源码、
\code{-O2} 优化级别和 10,000 次迭代的 NEMU 动态剖析进行 ISA 侧比较。未启用 B 扩展时执行
2,848,523,042 条指令；启用 Zba、Zbb、Zbc 和 Zbs 后执行 2,480,540,948 条，减少
367,982,094 条，即 12.918\%。后者包含 112,742,813 条 B 指令，占该镜像动态指令的 4.545\%。
这两个比例分别表示全程序动态指令缩减和 B 指令在新镜像中的占比，不能直接等同为处理器周期或板级
运行时间的缩减。

表~\ref{tab:b-hotspots} 列出历史全 B 镜像中出现次数最多的操作，并同时标出当前硬件是否保留。
\code{sh1add} 主要承担地址生成；
\code{sext.h}、\code{sh2add} 与 \code{pack/zext.h} 构成第二组热点。GNU 反汇编器会把
\code{pack rd,rs1,x0} 显示为 \code{zext.h}，因此二者合并计数。\code{clmul} 与 \code{clmulh}
各执行 5,840,008 次，合计 11,680,016 次，是 CRC 路径的关键替代操作。

\begin{table}[H]\centering\scriptsize
\caption{CoreMark 10,000 次迭代中的历史 B 指令热点及当前状态}\label{tab:b-hotspots}
\begin{tabularx}{\linewidth}{@{}lrrY@{}}
\toprule
\textbf{操作}&\textbf{动态次数}&\textbf{占全程序动态指令比例}&\textbf{提交 \texttt{8591368} 状态}\\\midrule
\code{sh1add}&60,280,009&2.4301\%&保留\\
\code{sext.h}&16,060,001&0.6474\%&保留\\
\code{sh2add}&12,120,153&0.4886\%&保留\\
\code{pack/zext.h}&10,261,412&0.4137\%&保留，由单周期 \code{pack} 实现\\
\code{clmul}&5,840,008&0.2354\%&移除；其已观测用途由 CRC 专用网络覆盖\\
\code{clmulh}&5,840,008&0.2354\%&移除；其已观测用途由 CRC 专用网络覆盖\\
\code{bext}&2,341,195&0.0944\%&保留\\
\bottomrule
\end{tabularx}
\end{table}

被移除的 \code{clmul/clmulh} 在历史镜像中合计占 0.4708\%，且本次观测到的热点用途集中在 CRC
更新；固定异或网络已覆盖该计算，同时避免保留通用迭代无进位乘法器。
其余被裁掉的迭代操作未进入上述热点集合；当前正式程序在迭代 B 单元中只使用
\code{clz/ctz}，因此删除它们可减少状态、选择和结果汇聚逻辑，不影响当前程序镜像。

1,000 次迭代的扩展子集剖析进一步给出：无 B 扩展为 284,879,140 条，仅启用 Zbc 为
251,874,680 条，启用 Zba、Zbb、Zbc、Zbs 全集为 248,080,124 条。由此可见，该工作负载的大部分
编译器侧收益来自 Zbc；这是历史 ISA 探索结果，不表示当前处理器仍支持 Zbc。全 B 镜像中的
\code{clmul/clmulh} 曾替代 CRC 热点内的大量移位和异或。
因此，在相同 B ISA 下启用 CRC 专用网络时，板测由 10.07521570~s 降至
9.56506582~s，严格同口径增量为 5.0634\%。历史上约 12.37\% 的数字混合了“无 B 软件 CRC”和
“有 B 且启用 CRC 加速”两个变化，不能归因于单一改动。受研发周期限制，本阶段未安排严格完整的
B-only 板级对照测试，因此本文不报告 B 扩展单独带来的板级百分比。

\subsection{专用硬件加速设计}\label{sec:custom}

动态剖析表明，热点主要集中在定宽位运算、局部乘累加、依赖访存和流式状态处理。设计分析分别识别出
地址生成、缓存查询、响应等待、局部计算和顺序提交等关键阶段，并针对各阶段配置专用组合或多周期硬件。
这些单元继续服从 EXU 的阻塞、冲刷与顺序写回协议，不改变输入数据、返回值或存储器可见顺序。

\begin{table}[H]\centering\scriptsize
\caption{专用硬件加速单元及其流水边界}\label{tab:accelerator-units}
\begin{tabularx}{\linewidth}{@{}p{0.19\linewidth}p{0.35\linewidth}p{0.19\linewidth}Y@{}}
\toprule
\textbf{硬件单元}&\textbf{主要结构}&\textbf{执行方式}&\textbf{体系结构边界}\\\midrule
CRC 字节更新&固定 XOR 网络与 16 位状态组合&单周期&无内部状态、无访存\\
窄位域乘法&位域抽取、小位宽乘法与补零&单周期&无内部状态、无访存\\
局部乘累加&窄乘法器与局部累加器&分阶段多周期&结果经长算术通路提交\\
访存型归约&地址生成与局部状态寄存器&阻塞式状态机&只读请求经统一 LSU 发出\\
依赖访存处理&D-cache 专用查询副本、统一请求口&阻塞式状态机&更新遵循普通 store 顺序\\
流式状态处理&对齐字读取与局部状态寄存器&有界多周期&只读输入，结果顺序提交\\
\bottomrule
\end{tabularx}
\normalsize\end{table}

\subsubsection{定宽位运算数据通路}

CRC 单元把一个输入字节与 16 位状态送入固定异或网络，直接得到下一状态；窄位域乘法单元并行完成
字段抽取和小位宽无符号乘法。两者均不保留跨操作状态，也不访问存储器，因此可接入单周期加速结果
通路，避免通用移位、掩码和异或序列反复占用整数数据通路。

\begin{figure}[H]\centering
\begin{tikzpicture}[n/.style={draw,rounded corners=1.2pt,minimum height=0.72cm,minimum width=2.15cm,align=center,fill=blue!6},a/.style={-Stealth,thick},font=\scriptsize,node distance=0.65cm and 0.85cm]
\node[n] (rs) {寄存器操作数\\rs1、rs2};
\node[n,above right=0.35cm and 1.1cm of rs] (crc) {8 位数据 + 16 位状态\\固定 XOR 矩阵};
\node[n,below right=0.35cm and 1.1cm of rs] (bmul) {位域抽取\\4 位 $\times$ 7 位};
\node[n,right=1.2cm of crc] (cz) {16 位 CRC\\高位补零};
\node[n,right=1.2cm of bmul] (bz) {11 位乘积\\高位补零};
\node[n,right=1.2cm of $(cz)!0.5!(bz)$] (wb) {硬件加速通路\\单周期写回};
\draw[a](rs)--(crc); \draw[a](rs)--(bmul); \draw[a](crc)--(cz); \draw[a](bmul)--(bz);
\draw[a](cz)--(wb); \draw[a](bz)--(wb);
\end{tikzpicture}
\caption{\texttt{xcrcu8} 与 \texttt{xbmul} 的定宽单周期数据通路}
\end{figure}


同 ISA 配置和迭代口径的板级对照中，CRC 硬件将运行时间由 10.07521570\,s 降至
9.56506582\,s，缩短 5.0634\%；继续加入窄位域乘法后降至 9.07727378\,s，相对前一点缩短
5.10\%。这两组结果用于说明定宽组合网络的收益，未将其外推到其他程序。

\subsubsection{乘累加与访存型归约}

乘累加单元复用窄 DSP，并在执行单元内部保存局部累加值。硬件按初始化、局部更新和结果整理等阶段
推进，减少中间数据在 GPR 与整数旁路上的往返；多周期握手保证相关消费者只在结果有效后继续执行。

访存型归约单元锁存请求参数后，由地址生成、请求、等待响应、局部变换和完成提交等状态逐拍推进。
每份输入数据仍经统一 LSU/D-cache 接口实际读取，单元不直接访问或旁路存储阵列。

以下四幅状态机图为便于说明控制关系，将多个硬件阶段合并画在同一视图中；
各节点对应独立的多拍步骤或握手边界，并非在一、两拍内直接产生整个操作的结果。

\begin{figure}[H]\centering
\resizebox{\linewidth}{!}{%
\begin{tikzpicture}[n/.style={draw,rounded corners=1.2pt,minimum height=0.78cm,minimum width=2.25cm,align=center,fill=green!7},a/.style={-Stealth,thick},font=\small,node distance=0.7cm]
\node[n] (cfg) {锁存 $base,n,clip$\\清零行列计数};
\node[n,right=of cfg] (req) {地址递增器\\发起 32 位读};
\node[n,right=of req] (resp) {响应寄存\\得到 $C_i$};
\node[n,right=of resp] (calc) {$T_i$、裁剪判断\\趋势计分};
\node[n,right=of calc] (next) {更新 $S_i,P_i,R_i$\\行列计数};
\node[n,right=of next] (done) {符号扩展\\写回 $R$};
\foreach \x/\y in {cfg/req,req/resp,resp/calc,calc/next,next/done}{\draw[a](\x)--(\y);}
\draw[a](next.south)--++(0,-0.45)-|node[pos=0.25,below]{尚有元素}(req.south);
\end{tikzpicture}}
\caption{\texttt{xmsum} 的顺序访存与递推归约流程}
\end{figure}


位域乘法融合版本在同一硬件配置下将板测时间由 6.690732940\,s 降至 6.485317060\,s。
加入局部乘累加状态后，相同 10/100 次 NPC 仿射模型的 10,000 次时间由 5.570051480\,s 降至
5.301518123\,s；该点未形成隔离板测，因此仅作为仿真拟合证据。访存型归约在对应板级对照中将
9.07727378\,s 降至 8.57272290\,s，缩短 5.56\%。

\subsubsection{依赖访存加速}

依赖访存的关键代价是后一次地址需要等待前一次数据。硬件在 D-cache 中设置与普通读取同步的专用查询
副本，并将并行查询、命中解析、统一总线回退、响应等待和完成提交划分为显式状态。设备区访问不会进入
缓存阵列。

\begin{figure}[H]\centering
\resizebox{\linewidth}{!}{%
\begin{tikzpicture}[
  s/.style={draw,rounded corners=1.2pt,align=center,minimum width=2.25cm,minimum height=0.78cm,fill=blue!6},
  d/.style={draw,diamond,aspect=2.15,align=center,fill=orange!8,inner sep=2pt},
  a/.style={-Stealth,thick},font=\scriptsize]
\node[s] (idle) at (6.0,6.5) {idle\\锁存参数};
\node[s] (lookup) at (6.0,5.2) {nodeLookup\\并行查 next/info};
\node[s] (resolve) at (6.0,3.9) {nodeResolve\\解析两路命中};
\node[s] (nextmem) at (2.3,3.0) {nextMemory\\等待 next};
\node[s] (inforeq) at (2.3,1.6) {infoRequest\\发出 info 请求};
\node[s] (infomem) at (6.0,1.6) {infoMemory\\等待 info};
\node[s] (dlookup) at (6.0,0.1) {dataLookup\\查询 $[info]$};
\node[s] (dresolve) at (9.6,0.1) {dataResolve\\解析数据命中};
\node[s] (dmem) at (9.6,-1.7) {dataMemory\\等待数据};
\node[d] (judge) at (6.0,-1.7) {匹配或到尾?};
\node[s] (done) at (2.3,-1.7) {done\\保持返回值};

\draw[a] (idle)--(lookup);
\draw[a] (lookup)--(resolve);
\draw[a] (resolve.west) -- ++(-0.55,0) -| node[pos=0.25,above]{next miss} (nextmem.north);
\draw[a] (resolve)--node[right]{next hit、info miss}(infomem);
\draw[a] (resolve.east) -- ++(0.55,0) -| node[pos=0.65,right]{两者命中} (dresolve.north);
\draw[a] (nextmem)--node[right]{info miss}(inforeq);
\draw[a] (nextmem.west) -- ++(-0.8,0) -- ++(0,-2.9) --
  node[above]{info hit} (dlookup.west);
\draw[a] (inforeq)--node[above]{请求握手}(infomem);
\draw[a] (infomem)--(dlookup);
\draw[a] (dlookup)--(dresolve);
\draw[a] (dresolve.south) to[out=245,in=35] node[pos=0.52,above,sloped]{数据命中} (judge.north);
\draw[a] (dresolve)--node[right]{数据 miss}(dmem);
\draw[a] (dmem)--node[above]{响应}(judge);
\draw[a] (judge)--node[above]{是}(done);
\draw[a] (judge.south) -- ++(0,-0.7) -- ++(5.3,0) |-
  node[pos=0.25,below]{否：装入 next} (lookup.east);
\draw[a] (done.west) -- ++(-1.5,0) |- node[pos=0.25,left]{写回握手} (idle.west);
\end{tikzpicture}}
\caption{\texttt{xlistfind} 的缓存命中/未命中状态机}\label{fig:custom-listfind}
\end{figure}


更新型单元同样以阻塞状态机保存局部状态。读数据经 D-cache 或统一总线取得，写入经普通 store 路径
提交；只有当前写入完成后才推进内部状态，因此保持原有的内存可见顺序。

\begin{figure}[H]\centering
\resizebox{\linewidth}{!}{%
\begin{tikzpicture}[n/.style={draw,rounded corners=1.2pt,minimum height=0.72cm,minimum width=1.8cm,align=center,fill=blue!6},d/.style={draw,diamond,aspect=2.0,align=center,fill=orange!8},a/.style={-Stealth,thick},font=\scriptsize,node distance=0.62cm and 0.68cm]
\node[n] (idle) {装入\\请求参数};
\node[d,right=of idle] (null) {可直接完成?};
\node[n,right=of null] (lookup) {查询 D-cache\\读取关联字段};
\node[d,right=of lookup] (hit) {命中?};
\node[n,right=of hit] (loadq) {发出 load};
\node[n,below=0.82cm of loadq] (loadr) {等待响应\\锁存数据};
\node[n,below=0.82cm of hit] (store) {发出 store\\提交局部更新};
\node[d,below=0.82cm of lookup] (more) {继续处理?};
\node[n,below=0.82cm of null] (done) {整理结果\\完成写回};
\draw[a](idle)--(null);
\draw[a](null)--node[above]{否}(lookup);
\draw[a](null)--node[right]{是}(done);
\draw[a](lookup)--(hit);
\draw[a](hit)--node[above]{否}(loadq);
\draw[a](hit)--node[right]{是}(store);
\draw[a](loadq)--(loadr); \draw[a](loadr)--(store); \draw[a](store)--(more);
\draw[a](more)--node[above]{否}(done);
\draw[a](more)--node[right]{是：更新状态}(lookup);
\end{tikzpicture}}
\caption{依赖访存更新单元的多拍控制}\label{fig:custom-listrev}
\end{figure}


在完整前序组合上加入依赖访存查询后，相同仿射模型的 10,000 次时间缩短 6.51\%，该点未形成隔离
板测。依赖访存更新的同口径
板级对照由 7.45481628\,s 降至 7.29058196\,s，缩短约 2.20\%。运行时间降幅小于指令数降幅，
与依赖访存的串行关系和存储访问延迟一致。

\subsubsection{流式状态处理}

流式状态单元把输入地址对齐到 32 位读取边界，并将响应整理、局部状态更新、继续处理和完成提交划分为
多个硬件阶段。该结构减少重复的字节加载和通用控制开销，但不改变输入缓冲区内容。

\begin{figure}[H]\centering
\resizebox{\linewidth}{!}{%
\begin{tikzpicture}[n/.style={draw,rounded corners=1.2pt,minimum height=0.82cm,minimum width=2.35cm,align=center,fill=purple!7},a/.style={-Stealth,thick},font=\small,node distance=0.68cm]
\node[n] (addr) {地址对齐\\发起字读取};
\node[n,right=of addr] (reg) {响应寄存\\整理有效字节};
\node[n,right=of reg] (lo) {状态更新阶段 0\\处理局部输入};
\node[n,right=of lo] (hi) {状态更新阶段 1\\延续局部状态};
\node[n,right=of hi] (pack) {保存局部状态\\整理返回信息};
\node[n,right=of pack] (commit) {满足完成条件\\提交结果};
\foreach \x/\y in {addr/reg,reg/lo,lo/hi,hi/pack,pack/commit}{\draw[a](\x)--(\y);}
\draw[a](pack.south)--++(0,-0.45)-|node[pos=0.25,below]{仍有输入：继续处理}(addr.south);
\end{tikzpicture}}
\caption{流式状态单元的多拍处理阶段}
\end{figure}


流式状态硬件加入对应组合后，板测时间由 6.485317060\,s 降至 6.026062020\,s，缩短 7.081\%。
另一版组合板测时间为 6.69637186\,s；该结果只作为组合证据，不归因于单一单元。

\subsubsection{加速构建与验证口径}

程序镜像按目标硬件能力配置构建；启用与关闭版本采用相同输入、迭代规模和结果判定，并分别完成差分与
CRC 验证。
最终加速构建在 NPC 中以 10/100 次结果拟合得到
\textbf{1,587,738,705 cycles，即 5.292462350\,s@300\,MHz}，四项 CRC 均正确
\evidence{E-ACCEL-FINAL}；精确 NEMU 10,000 次运行提交 773,441,484 条指令并通过相同 CRC。
该时间尚未上板。构建配置调整未改变 RTL，因此时序结果沿用同一硬件实现。
\subsection{一致性、阻塞与异常边界}

无内部访存的定宽位运算进入加速结果通路；局部乘累加进入复用窄 DSP 的长算术通路；其余单元在 EXU
保持当前操作，直至状态机完成并与后级握手。依赖访存和流式状态单元的读取均复用统一数据总线。
控制流冲刷不会提交年轻操作的请求；已经握手的写入按普通 store 的可见顺序完成。每个单元只驱动唯一 ResultKind lane，使阻塞控制、访存副作用
和普通整数结果选择保持清晰边界。
